
\documentclass[12pt,a4paper]{article}
% Compilacao: pdflatex relatorio.tex (duas passagens)
\usepackage[utf8]{inputenc}
\usepackage[T1]{fontenc}
\usepackage[brazil]{babel}
\usepackage{lmodern}
\usepackage{geometry,setspace,amsmath,amssymb,booktabs,array,longtable}
\usepackage{graphicx,float,tikz,pgfplots,listings,enumitem}
\usepackage{fancyhdr,hyperref,microtype}
\pgfplotsset{compat=1.18}
\geometry{a4paper,left=3cm,right=2cm,top=2.5cm,bottom=2.5cm}
\onehalfspacing
\setlength{\parindent}{1.25cm}
\setlength{\parskip}{0.2em}
\setlength{\headheight}{15pt}
\pagestyle{fancy}
\fancyhf{}
\fancyhead[R]{\small Problema dos $k$-centros}
\fancyfoot[C]{\thepage}
\renewcommand{\headrulewidth}{0.4pt}
\hypersetup{colorlinks=true,linkcolor=black,urlcolor=black,
  pdftitle={Problema dos k-centros: metodos exato e aproximado},
  pdfauthor={Rafael Xavier Oliveira; Thales Duque Camara; Gabriel Teotonio de Castro Coelho Costa}}
\lstset{basicstyle=\ttfamily\footnotesize,columns=fullflexible,
  breaklines=true,showstringspaces=false,frame=single,
  numbers=none,aboveskip=6pt,belowskip=6pt}
\setlist[itemize]{itemsep=3pt,topsep=5pt}
\newcommand{\bigO}[1]{\mathcal O\!\left(#1\right)}
\begin{document}
\begin{titlepage}
\begin{center}
{\large Pontifícia Universidade Católica de Minas Gerais}\\[0.7cm]
{\large Unidade Lourdes --- Ciência da Computação}\\[3.6cm]
{\large Teoria dos Grafos e Computabilidade}\\[0.4cm]
{\large Trabalho Prático n.º 02}\\[1cm]
{\bfseries\LARGE O PROBLEMA DOS $k$-CENTROS}\\[0.6cm]
{\large Implementação e comparação de uma solução\\
exata e uma solução aproximada}\\[2.8cm]
{\large Rafael Xavier Oliveira}\\[0.28cm]
{\large Thales Duque Câmara}\\[0.28cm]
{\large Gabriel Teotônio de Castro Coelho Costa}\\[0.75cm]
{\large Professor: Zenilton Kleber Gonçalves do Patrocínio Júnior}\\[2.0cm]
{\large Belo Horizonte --- 2026}
\end{center}
\end{titlepage}

\section{Introdução}
O problema dos $k$-centros consiste em selecionar até $k$ vértices de um
conjunto de locais para minimizar a distância máxima que qualquer ponto
precisará percorrer até o centro selecionado mais próximo. O critério é de
\emph{minimização do pior atendimento}: diferentemente do problema das
$p$-medianas, o objetivo não é minimizar a soma de todas as distâncias.
Essa distinção é essencial para interpretar os dados e os valores de referência.

O enunciado do Trabalho Prático n.º 02 solicita duas formas de resolução:
um método exato, capaz de encontrar e demonstrar o valor ótimo em
instâncias pequenas, e outro aproximado, aplicável a instâncias de maior porte.
A implementação foi feita em Java, com leitura do formato OR-Library,
preparação da matriz de distâncias, testes automatizados e exportação dos
resultados em arquivos CSV. Conservou-se, como referência de organização,
a estrutura em pastas e a interface por linha de comando do Trabalho Prático
n.º 01, mas os algoritmos são específicos do novo problema.

\section{Definição do problema}
Seja $V$ um conjunto de $n$ vértices e $d(u,v)$ a menor distância entre
$u$ e $v$. Deseja-se escolher $C\subseteq V$, com $1\leq |C|\leq k$, minimizando
\begin{equation}\label{eq:objetivo}
 r(C)=\max_{v\in V}\min_{c\in C}d(v,c).
\end{equation}
O valor $r(C)$ é chamado \emph{raio}. Para cada vértice, considera-se o
centro mais próximo; entre os vértices, considera-se a maior dessas distâncias.
O raio ótimo é $r^*=\min_{C,\,|C|\leq k}r(C)$.

A matriz de distâncias deve satisfazer simetria, não negatividade,
$d(v,v)=0$ e desigualdade triangular. Essas propriedades passam a valer
ao substituir os pesos das arestas originais por distâncias de caminhos
mínimos em um grafo não direcionado com custos não negativos.

\subsection{Preparação das instâncias OR-Library}
Cada arquivo \texttt{pmedN.txt} contém, na primeira linha, $n$, $m$ e $k$.
Nas $m$ linhas seguintes aparecem vértice inicial, vértice final e custo
da aresta. Os identificadores de vértices no arquivo começam em 1;
internamente, passam a começar em 0. Os dados descrevem um \emph{grafo
esparso}, e não uma matriz completa de distâncias. Por isso, antes de
resolver $k$-centros, aplica-se Floyd--Warshall:
\begin{equation}
 d^{(t)}(i,j)=\min\left\{d^{(t-1)}(i,j),\,
 d^{(t-1)}(i,t)+d^{(t-1)}(t,j)\right\}.
\end{equation}
A página oficial da OR-Library especifica um detalhe importante: se uma
aresta não direcionada aparecer mais de uma vez, \textbf{o último custo
lido substitui os anteriores}. A implementação segue exatamente essa regra,
em vez de escolher indevidamente o menor custo entre duplicatas.
O pré-processamento exige $\bigO{n^3}$ tempo e $\bigO{n^2}$ memória.

\section{Algoritmo aproximado}
Foi adotada a estratégia \emph{farthest-first} de Gonzalez, com início
arbitrário: escolhe-se o primeiro centro, depois, a cada passo, seleciona-se
o vértice cuja distância ao centro mais próximo é máxima. A cada inserção,
atualiza-se para todos os vértices a menor distância a um centro escolhido.
A técnica tem garantia de fator 2 em espaços métricos:
\begin{equation}\label{eq:approx}
 r_{\mathrm{FF}}\leq 2r^*.
\end{equation}

\begin{lstlisting}[caption={Construção aproximada dos centros (uma inicialização).}]
GONZALEZ(d, k, inicio):
    centros = [inicio]
    para cada vertice v: perto[v] = d(v, inicio)
    enquanto tamanho(centros) < k:
        c = vertice nao escolhido com maior perto[c]
        adicionar c a centros
        para cada vertice v:
            perto[v] = min(perto[v], d(v, c))
    retornar centros
\end{lstlisting}

\noindent\textbf{Aperfeiçoamentos implementados.}
Em vez de usar só um ponto inicial, executa-se a construção a partir de
quatro posições determinísticas ($0$, $\lfloor n/3\rfloor$,
$\lfloor 2n/3\rfloor$ e $n-1$) e escolhe-se a menor solução.
Depois, tenta-se uma busca local com até duas passagens de trocas entre
um centro ativo e um vértice não selecionado. Uma troca só é aceita se
reduzir estritamente o raio. Usando as distâncias para o primeiro e para
o segundo centro mais próximos, cada avaliação de troca requer
$\bigO{n}$. O refinamento não invalida a garantia da
Equação~\ref{eq:approx}, pois nunca aumenta o raio.

\noindent\textbf{Custo.}
Cada construção custa $\bigO{nk}$; quatro construções preservam essa
ordem. Com $t$ passagens da busca local, o limite superior é
$\bigO{tn^2k}$, além do pré-processamento comum.

\subsection{Justificativa do fator 2}
Ao selecionar $k$ centros pelo método mais distante e considerar um
vértice adicional a distância $R$ da seleção, obtêm-se $k+1$ pontos
separados dois a dois por pelo menos $R$. Como uma solução ótima distribui
esses pontos entre apenas $k$ centros, dois necessariamente pertencem à
mesma região atendida. Pela desigualdade triangular, a distância entre
esses dois pontos é no máximo $2r^*$. Portanto, $R\leq 2r^*$.
A escolha do melhor entre vários inícios e a melhoria por trocas só
reduzem o raio, mantendo a mesma cota.

\section{Algoritmo exato}
O método exato transforma a busca em um conjunto finito de problemas
de \emph{cobertura} e os resolve com busca em profundidade e
\emph{branch-and-bound}. Para um raio candidato $R$, cada vértice $c$
representa o conjunto de vértices atendidos
\begin{equation}
 S_c(R)=\{v\in V : d(v,c)\leq R\}.
\end{equation}
Existe uma solução de raio no máximo $R$ exatamente quando se pode
selecionar até $k$ desses conjuntos de forma que a união cubra todo $V$.
A decisão equivale ao sistema binário:
\begin{align}
 \sum_{c:\,d(v,c)\leq R}x_c &\geq 1 &&\forall\,v\in V,\\
 \sum_{c\in V}x_c &\leq k,\\
 x_c &\in\{0,1\}.\nonumber
\end{align}

Os raios candidatos são os valores distintos da matriz de distâncias.
Uma busca binária verifica o menor raio viável, começando pelo intervalo
limitado inferiormente por $\lceil r_{\mathrm{aprox}}/2\rceil$, consequência
da garantia de fator 2, e superiormente pelo raio aproximado.
O algoritmo de decisão utiliza vetores de bits \texttt{BitSet} para armazenar
coberturas. A ramificação escolhe um vértice ainda não atendido e testa
possíveis centros que o cobrem, priorizando centros que atendem mais pontos.

\begin{lstlisting}[caption={Decisao exata para raio R.}]
COBRE(R, centros_restantes, descobertos):
    se todos os vertices estiverem cobertos: retornar VIAVEL
    se centros_restantes == 0: retornar INVIAVEL
    se teto(|descobertos| / maior_cobertura) > centros_restantes:
        retornar INVIAVEL
    se existem mais de centros_restantes vertices
       mutuamente separados por distancias > 2*R:
        retornar INVIAVEL
    v = vertice descoberto com menos candidatos de centro
    para cada candidato c que cobre v (maior ganho primeiro):
        se COBRE(R, centros_restantes-1,
                 descobertos sem cobertura(c)) == VIAVEL:
            registrar c; retornar VIAVEL
    retornar INVIAVEL
\end{lstlisting}

Também se eliminam ramificações com conjunto descoberto idêntico.
A condição de empacotamento é válida porque dois vértices separados por
mais de $2R$ não podem ser atendidos pelo mesmo centro. Quando a solução
usa menos de $k$ centros, centros adicionais podem ser incluídos sem
prejudicar o raio, de forma que a saída tenha $k$ centros.

\noindent\textbf{Critério de parada.}
A busca exata possui pior caso exponencial (a decisão é uma variante
de cobertura de conjuntos). Por isso, há um limite configurável de tempo
por instância. Se o prazo termina durante uma decisão, o programa retorna
a melhor solução válida conhecida e informa \texttt{TEMPO\_ESGOTADO};
\emph{nesse caso não afirma otimalidade}. O status
\texttt{OPTIMO\_COMPROVADO} é emitido somente quando as decisões necessárias
foram concluídas, provando a inexistência de raio melhor.






























\section{Implementação e validação}
A solução foi implementada em Java e organizada conforme o TP01 da turma:
\texttt{grafos/} para os fontes, \texttt{instancias/} para as entradas,
\texttt{resultados/} para as evidências e \texttt{relatorio/} para TeX e PDF.
As 40 entradas oficiais estão incluídas no pacote, bem como dez instâncias
sintéticas adicionais. Cada arquivo oficial foi conferido por quantidade de
vértices, arestas, centros e assinatura numérica (FNV-1a).

\begin{table}[H]\centering\small
\caption{Componentes da implementação.}
\begin{tabular}{@{}p{5cm}p{7.2cm}@{}}\toprule
Arquivo & Função \\ \midrule
\texttt{GraphIO.java} & Leitura das instâncias OR-Library, tratamento de duplicatas e caminhos mínimos por Floyd--Warshall.\\
\texttt{ApproxKCenter.java} & Escolha do vértice mais distante, quatro sementes e duas passagens de trocas locais.\\
\texttt{ExactKCenter.java} & Busca binária no raio, decisão por cobertura e branch-and-bound.\\
\texttt{Benchmark.java} & Execução dos dois métodos, limites de tempo e medição por instância.\\
\texttt{GraphGenerator.java} & Preparação dos grafos sintéticos usados nos testes adicionais.\\
\texttt{SelfTest.java} & Testes contra enumeração completa em casos pequenos.\\
\bottomrule\end{tabular}
\end{table}

A suíte automatizada aprovou \textbf{131 verificações}, incluindo tratamento de
arestas repetidas, caminhos mínimos, garantia de aproximação e igualdade
entre algoritmo exato e enumeração exaustiva em exemplos pequenos.
Como verificação adicional, as 80 escolhas de centros (40 do método aproximado
e 40 da busca exata) foram conferidas por uma rotina independente em Python,
com distâncias recalculadas por Dijkstra multiorigem (Python); os raios registrados em Java
coincidiram com os raios calculados independentemente. Essa auditoria confirma
\emph{validade dos centros e dos raios}; ela não prova que uma solução
interrompida seja ótima.

\section{Metodologia experimental}
\subsection{Dados e métricas}
Foram utilizadas \textbf{todas as 40 instâncias} \texttt{pmed1.txt} a
\texttt{pmed40.txt} da OR-Library, de 100 a 900 vértices, com os valores
$k$ e os raios de referência $r^*$ fornecidos na Tabela 1 do enunciado.
As instâncias da OR-Library são inicialmente grafos ponderados esparsos;
a matriz completa de distância foi produzida por Floyd--Warshall.
Arestas paralelas foram tratadas de acordo com a especificação das entradas:
a última ocorrência define o peso armazenado. Não foram usados os ótimos
de $p$-medianas como referência para a função objetivo de $k$-centros.

Para medir a eficácia, utilizou-se o raio $r(C)$ de cada solução, a diferença
em relação ao valor de referência e o percentual
\[
\operatorname{gap}(C)=100\frac{r(C)-r^*}{r^*}\,\%.
\]
A eficiência foi medida em milissegundos com \texttt{System.nanoTime()},
separadamente para pré-processamento, heurística e busca exata.
O orçamento por instância do exato foi de \textbf{3000~ms}; na avaliação
sintética, de 1000~ms. Resultados \texttt{TEMPO\_ESGOTADO} indicam
soluções viáveis, \emph{não} certificados de otimalidade. Tempos
com corte não representam o tempo necessário para terminar o algoritmo.
Os valores são de uma execução por entrada, sem repetições ou intervalos
de confiança; devem ser interpretados como estudo experimental descritivo.

\section{Resultados das 40 instâncias oficiais}
\subsection{Qualidade das soluções}
\begin{longtable}{@{}lrrrrrrc@{}}
\caption{Resultados Java: raios das duas abordagens nas 40 instâncias; estado do método exato.}\\
\toprule Caso & $n$ & $k$ & $r^*$ & $r_A$ & Gap (\%) & $r_E$ & Estado\\ \midrule
\endfirsthead
\toprule Caso & $n$ & $k$ & $r^*$ & $r_A$ & Gap (\%) & $r_E$ & Estado\\ \midrule
\endhead
pmed1 & 100 & 5 & 127 & 149 & 17,32 & 127 & OC \\
pmed2 & 100 & 10 & 98 & 117 & 19,39 & 98 & OC \\
pmed3 & 100 & 10 & 93 & 123 & 32,26 & 93 & OC \\
pmed4 & 100 & 20 & 74 & 98 & 32,43 & 74 & OC \\
pmed5 & 100 & 33 & 48 & 64 & 33,33 & 48 & OC \\
pmed6 & 200 & 5 & 84 & 96 & 14,29 & 84 & OC \\
pmed7 & 200 & 10 & 64 & 86 & 34,38 & 64 & TE \\
pmed8 & 200 & 20 & 55 & 71 & 29,09 & 71 & TE \\
pmed9 & 200 & 40 & 37 & 51 & 37,84 & 38 & TE \\
pmed10 & 200 & 67 & 20 & 29 & 45,00 & 20 & TE \\
pmed11 & 300 & 5 & 59 & 61 & 3,39 & 59 & OC \\
pmed12 & 300 & 10 & 51 & 65 & 27,45 & 65 & TE \\
pmed13 & 300 & 30 & 35 & 53 & 51,43 & 40 & TE \\
pmed14 & 300 & 60 & 26 & 36 & 38,46 & 27 & TE \\
pmed15 & 300 & 100 & 18 & 24 & 33,33 & 18 & TE \\
pmed16 & 400 & 5 & 47 & 51 & 8,51 & 47 & OC \\
pmed17 & 400 & 10 & 39 & 47 & 20,51 & 47 & TE \\
pmed18 & 400 & 40 & 28 & 41 & 46,43 & 31 & TE \\
pmed19 & 400 & 80 & 18 & 28 & 55,56 & 21 & TE \\
pmed20 & 400 & 133 & 13 & 18 & 38,46 & 13 & TE \\
pmed21 & 500 & 5 & 40 & 44 & 10,00 & 40 & OC \\
pmed22 & 500 & 10 & 38 & 49 & 28,95 & 49 & TE \\
pmed23 & 500 & 50 & 22 & 32 & 45,45 & 24 & TE \\
pmed24 & 500 & 100 & 15 & 22 & 46,67 & 16 & TE \\
pmed25 & 500 & 167 & 11 & 15 & 36,36 & 11 & TE \\
pmed26 & 600 & 5 & 38 & 42 & 10,53 & 38 & OC \\
pmed27 & 600 & 10 & 32 & 37 & 15,62 & 33 & TE \\
pmed28 & 600 & 60 & 18 & 28 & 55,56 & 21 & TE \\
pmed29 & 600 & 120 & 13 & 19 & 46,15 & 14 & TE \\
pmed30 & 600 & 200 & 9 & 13 & 44,44 & 10 & TE \\
pmed31 & 700 & 5 & 30 & 35 & 16,67 & 30 & OC \\
pmed32 & 700 & 10 & 29 & 35 & 20,69 & 35 & TE \\
pmed33 & 700 & 70 & 15 & 25 & 66,67 & 19 & TE \\
pmed34 & 700 & 140 & 11 & 17 & 54,55 & 13 & TE \\
pmed35 & 800 & 5 & 30 & 34 & 13,33 & 30 & OC \\
pmed36 & 800 & 10 & 27 & 34 & 25,93 & 34 & TE \\
pmed37 & 800 & 80 & 15 & 24 & 60,00 & 18 & TE \\
pmed38 & 900 & 5 & 29 & 31 & 6,90 & 29 & OC \\
pmed39 & 900 & 10 & 23 & 30 & 30,43 & 30 & TE \\
pmed40 & 900 & 90 & 13 & 20 & 53,85 & 15 & TE \\
\bottomrule
\end{longtable}
\noindent\textbf{Legenda:} OC = ótimo comprovado; TE = tempo esgotado,
com melhor solução viável. $r_A$ = raio aproximado; $r_E$ = melhor raio
encontrado pela busca exata; $r^*$ = referência da atividade.

A solução aproximada teve gap médio de \textbf{32,69\%}
(mediana de 32,88\%), com mínimo de 3,39\% em
\texttt{pmed11} e máximo de 66,67\% em \texttt{pmed33}.
A busca exata reduziu estritamente o raio da aproximação em
\textbf{33 dos 40 casos}; o gap médio do melhor incumbente foi de
\textbf{9,15\%}. Em 18 casos, esse raio coincidiu com
a referência, mas apenas 13 obtiveram certificação durante a busca:
os 27 restantes atingiram o limite de 3000~ms.

\begin{figure}[H]\centering
\begin{tikzpicture}
\begin{axis}[width=0.94\textwidth,height=5.4cm,xmin=0,xmax=41,
 xtick={1,5,10,15,20,25,30,35,40},ylabel={Gap da heurística (\%)},
 xlabel={Índice da instância},ymajorgrids=true,grid style={gray!20}]
\addplot[only marks,mark=*,black,mark size=1.8pt] coordinates {(1,17.3228) (2,19.3878) (3,32.2581) (4,32.4324) (5,33.3333) (6,14.2857) (7,34.3750) (8,29.0909) (9,37.8378) (10,45.0000) (11,3.3898) (12,27.4510) (13,51.4286) (14,38.4615) (15,33.3333) (16,8.5106) (17,20.5128) (18,46.4286) (19,55.5556) (20,38.4615) (21,10.0000) (22,28.9474) (23,45.4545) (24,46.6667) (25,36.3636) (26,10.5263) (27,15.6250) (28,55.5556) (29,46.1538) (30,44.4444) (31,16.6667) (32,20.6897) (33,66.6667) (34,54.5455) (35,13.3333) (36,25.9259) (37,60.0000) (38,6.8966) (39,30.4348) (40,53.8462)};
\end{axis}\end{tikzpicture}
\caption{Desvio percentual do método aproximado nas 40 instâncias oficiais.}
\end{figure}

\subsection{Tempos observados e limites da busca}
\begin{longtable}{@{}lrrrr@{}}
\caption{Tempos dos experimentos Java em milissegundos (uma execução por entrada).}\\
\toprule Caso & $t_{FW}$ (ms) & $t_A$ (ms) & $t_E$ (ms) & LB do exato\\ \midrule
\endfirsthead
\toprule Caso & $t_{FW}$ (ms) & $t_A$ (ms) & $t_E$ (ms) & LB do exato\\ \midrule
\endhead
pmed1 & 4,92 & 2,79 & 123,73 & 127 \\
pmed2 & 6,03 & 4,65 & 564,50 & 98 \\
pmed3 & 5,58 & 6,41 & 281,49 & 93 \\
pmed4 & 7,39 & 8,02 & 1917,06 & 74 \\
pmed5 & 5,61 & 8,30 & 354,59 & 48 \\
pmed6 & 29,67 & 6,49 & 408,84 & 84 \\
pmed7 & 32,18 & 9,99 & 3001,40 & 59 \\
pmed8 & 44,19 & 14,07 & 3119,99 & 36 \\
pmed9 & 63,21 & 32,16 & 3002,09 & 32 \\
pmed10 & 28,51 & 26,71 & 3001,77 & 19 \\
pmed11 & 158,43 & 19,53 & 149,19 & 59 \\
pmed12 & 164,98 & 17,06 & 3002,51 & 33 \\
pmed13 & 155,69 & 34,30 & 3002,65 & 27 \\
pmed14 & 201,76 & 38,95 & 3050,10 & 23 \\
pmed15 & 84,88 & 34,76 & 3002,88 & 15 \\
pmed16 & 145,18 & 20,29 & 136,52 & 47 \\
pmed17 & 142,91 & 9,62 & 3001,14 & 24 \\
pmed18 & 147,43 & 46,50 & 3003,86 & 21 \\
pmed19 & 141,24 & 42,10 & 3001,26 & 14 \\
pmed20 & 134,26 & 47,73 & 3001,10 & 11 \\
pmed21 & 294,85 & 14,20 & 1098,90 & 40 \\
pmed22 & 426,19 & 37,45 & 3002,34 & 25 \\
pmed23 & 347,73 & 118,46 & 3001,07 & 16 \\
pmed24 & 239,09 & 133,99 & 3001,91 & 11 \\
pmed25 & 273,93 & 115,61 & 3001,43 & 8 \\
pmed26 & 429,81 & 30,87 & 786,14 & 38 \\
pmed27 & 336,14 & 68,90 & 3001,13 & 29 \\
pmed28 & 476,89 & 64,42 & 3004,60 & 14 \\
pmed29 & 1081,35 & 187,47 & 3001,84 & 10 \\
pmed30 & 538,20 & 110,69 & 3002,54 & 7 \\
pmed31 & 835,02 & 19,64 & 2790,40 & 30 \\
pmed32 & 526,99 & 41,93 & 3001,42 & 18 \\
pmed33 & 520,90 & 86,60 & 3002,81 & 13 \\
pmed34 & 513,50 & 111,02 & 3000,89 & 9 \\
pmed35 & 1034,18 & 31,56 & 557,17 & 30 \\
pmed36 & 1331,76 & 87,78 & 3003,23 & 17 \\
pmed37 & 792,11 & 115,53 & 3001,57 & 12 \\
pmed38 & 1191,76 & 23,60 & 225,54 & 29 \\
pmed39 & 1255,04 & 49,03 & 3002,51 & 15 \\
pmed40 & 1207,54 & 307,02 & 3002,38 & 10 \\
\bottomrule
\end{longtable}
O campo LB registra o limite inferior retornado pelo método exato. Se
a busca termina, o valor coincide com o ótimo. Nos casos TE,
é apenas uma cota inferior conhecida, com solução viável acima ou igual.

\begin{figure}[H]\centering
\begin{tikzpicture}
\begin{semilogyaxis}[width=0.94\textwidth,height=5.6cm,
 xmin=0,xmax=41,xtick={1,5,10,15,20,25,30,35,40},
 xlabel={Índice da instância},ylabel={Tempo (ms; escala logarítmica)},
 ymajorgrids=true,grid style={gray!20},legend pos=north west]
\addplot[only marks,mark=*,black] coordinates {(1,2.7870) (2,4.6520) (3,6.4050) (4,8.0180) (5,8.3020) (6,6.4930) (7,9.9910) (8,14.0670) (9,32.1590) (10,26.7060) (11,19.5340) (12,17.0630) (13,34.2980) (14,38.9530) (15,34.7600) (16,20.2850) (17,9.6240) (18,46.4990) (19,42.0970) (20,47.7320) (21,14.1990) (22,37.4480) (23,118.4580) (24,133.9900) (25,115.6150) (26,30.8730) (27,68.8960) (28,64.4230) (29,187.4710) (30,110.6900) (31,19.6370) (32,41.9320) (33,86.6030) (34,111.0220) (35,31.5580) (36,87.7770) (37,115.5290) (38,23.6000) (39,49.0310) (40,307.0190)};
\addplot[only marks,mark=square*,gray] coordinates {(1,123.7270) (2,564.5000) (3,281.4910) (4,1917.0620) (5,354.5890) (6,408.8390) (7,3001.3980) (8,3119.9880) (9,3002.0870) (10,3001.7710) (11,149.1920) (12,3002.5070) (13,3002.6510) (14,3050.1020) (15,3002.8780) (16,136.5180) (17,3001.1360) (18,3003.8600) (19,3001.2570) (20,3001.0970) (21,1098.9040) (22,3002.3450) (23,3001.0710) (24,3001.9060) (25,3001.4250) (26,786.1450) (27,3001.1320) (28,3004.5980) (29,3001.8400) (30,3002.5400) (31,2790.4040) (32,3001.4230) (33,3002.8130) (34,3000.8910) (35,557.1660) (36,3003.2340) (37,3001.5670) (38,225.5360) (39,3002.5080) (40,3002.3840)};
\legend{Aproximado,Exato}
\end{semilogyaxis}\end{tikzpicture}
\caption{Tempos dos dois métodos; valores do exato incluem encerramentos por limite.}
\end{figure}

\subsection{Resultados agrupados por quantidade de vértices}
\begin{table}[H]\centering\scriptsize
\caption{Gaps médios por tamanho, medianas de tempo e certificados da busca exata.}
\resizebox{\textwidth}{!}{\begin{tabular}{@{}rrrrrrr@{}}\toprule
$n$ & Casos & Gap aprox. (\%) & Gap exato (\%) & Mediana $t_A$ & Mediana $t_E$ & Provas\\ \midrule
100 & 5 & 26,95 & 0,00 & 6,41 & 354,59 & 5 \\
200 & 5 & 32,12 & 6,36 & 14,07 & 3001,77 & 1 \\
300 & 5 & 30,81 & 9,12 & 34,30 & 3002,65 & 1 \\
400 & 5 & 33,89 & 9,58 & 42,10 & 3001,14 & 1 \\
500 & 5 & 33,49 & 8,94 & 115,61 & 3001,43 & 1 \\
600 & 5 & 34,46 & 7,72 & 68,90 & 3001,84 & 1 \\
700 & 4 & 39,64 & 16,38 & 64,27 & 3001,16 & 1 \\
800 & 3 & 33,09 & 15,31 & 87,78 & 3001,57 & 1 \\
900 & 3 & 30,39 & 15,27 & 49,03 & 3002,38 & 1 \\
\bottomrule\end{tabular}}
\end{table}
\begin{figure}[H]\centering
\begin{tikzpicture}
\begin{axis}[width=0.94\textwidth,height=5.4cm,xlabel={Número de vértices},
 ylabel={Gap médio (\%)},xtick={100,200,300,400,500,600,700,800,900},
 ymajorgrids=true,grid style={gray!20},legend pos=north west]
\addplot[black,mark=*] coordinates {(100,26.9469) (200,32.1179) (300,30.8129) (400,33.8938) (500,33.4864) (600,34.4610) (700,39.6421) (800,33.0864) (900,30.3925)};
\addplot[gray,mark=square*] coordinates {(100,0.0000) (200,6.3587) (300,9.1166) (400,9.5788) (500,8.9410) (600,7.7190) (700,16.3845) (800,15.3086) (900,15.2731)};
\legend{Aproximado,Busca exata (melhor raio)}
\end{axis}\end{tikzpicture}
\caption{Qualidade média por quantidade de vértices, sem isolar o efeito de $k$.}
\end{figure}

A mediana de tempo da aproximação foi de \textbf{34,53~ms}; a
busca exata teve mediana de \textbf{3001,41~ms}, influenciada
pelo limite imposto. O pré-processamento por Floyd--Warshall levou, em média,
383,93~ms (máximo de 1331,76~ms), custo comum às
duas abordagens. Entre grupos, tanto o parâmetro $k$ quanto a estrutura
das distâncias variam. Assim, diferenças de gap entre tamanhos não
podem ser atribuídas exclusivamente ao número de vértices.

\section{Experimentos complementares em grafos sintéticos}
O conjunto adicional compreende dez grafos gerados por sementes fixas,
com tamanhos de 10 a 900 vértices. São úteis para observar comportamento,
mas \emph{não substituem} as 40 instâncias oficiais obrigatórias.
\begin{table}[H]\centering\scriptsize
\caption{Resultados sintéticos Java; tempos em milissegundos.}
\resizebox{\textwidth}{!}{\begin{tabular}{@{}rrrrrrrc@{}}\toprule
$n$ & $k$ & $t_{FW}$ & $r_A$ & $t_A$ & $r_E$ & $t_E$ & Estado\\ \midrule
10 & 2 & 0,03 & 85 & 0,98 & 85 & 2,14 & OC \\
12 & 3 & 0,06 & 49 & 0,37 & 45 & 0,34 & OC \\
16 & 3 & 0,18 & 75 & 0,15 & 75 & 0,42 & OC \\
20 & 4 & 0,30 & 86 & 0,10 & 86 & 0,84 & OC \\
40 & 5 & 2,23 & 94 & 0,66 & 86 & 10,89 & OC \\
100 & 10 & 8,96 & 117 & 2,98 & 90 & 856,23 & OC \\
200 & 20 & 17,48 & 126 & 15,64 & 126 & 1000,28 & TE \\
400 & 40 & 44,46 & 138 & 113,16 & 103 & 1000,21 & TE \\
600 & 60 & 264,24 & 146 & 24,10 & 109 & 1000,18 & TE \\
900 & 90 & 389,52 & 141 & 108,68 & 106 & 1000,17 & TE \\
\bottomrule\end{tabular}}
\end{table}
\begin{figure}[H]\centering
\begin{tikzpicture}
\begin{semilogyaxis}[width=0.94\textwidth,height=5.2cm,
 xmin=0,xmax=950,xtick={10,100,200,400,600,900},
 xlabel={Número de vértices},ylabel={Tempo (ms, escala logarítmica)},
 ymajorgrids=true,grid style={gray!20},legend pos=north west]
\addplot[black,mark=*] coordinates {(10,0.9840) (12,0.3670) (16,0.1510) (20,0.0980) (40,0.6610) (100,2.9750) (200,15.6380) (400,113.1630) (600,24.0980) (900,108.6840)};
\addplot[gray,mark=square*] coordinates {(10,0.0320) (12,0.0600) (16,0.1780) (20,0.3000) (40,2.2320) (100,8.9620) (200,17.4770) (400,44.4570) (600,264.2380) (900,389.5250)};
\legend{Aproximação,Floyd--Warshall}
\end{semilogyaxis}\end{tikzpicture}
\caption{Escalabilidade da aproximação e preparação no conjunto sintético.}
\end{figure}

\section{Discussão dos resultados}
A heurística de Gonzalez fornece uma garantia de fator 2 para distâncias
métricas. As sementes múltiplas e as trocas locais não pioram o raio,
portanto mantêm essa garantia. Sua grande vantagem é o custo computacional
polinomial, que permitiu resolver todas as entradas oficiais rapidamente
após o pré-processamento comum.

O método exato emprega teste de cobertura e explora possibilidades de
selecionar centros. A combinação de busca binária no raio, limitação por
cobertura máxima e empacotamento reduz muitos ramos, mas a complexidade
combinatória permanece. A comparação evidencia essa diferença: o método
exato conseguiu melhorar a incumbente em 33 casos, porém não
certificou o ótimo em 27 instâncias ao orçamento escolhido.
Quando o valor de referência é alcançado num caso TE, isso não altera
o status: não houve certificação pela busca implementada.

É importante distinguir \emph{eficácia} (qualidade do raio) de
\emph{eficiência} (tempo). O exato exibiu frequentemente um raio menor,
mas consumiu muito mais tempo; o aproximado foi mais rápido e teve
uma garantia teórica mesmo quando seu raio ficou acima da referência.
As medições de uma única execução e a influência do JIT na JVM impedem
uma conclusão estatística sobre tempos precisos. A análise é comparativa
e descritiva, apoiada por todas as 40 instâncias oficiais.

\section{Reprodutibilidade}
Os 40 arquivos oficiais da OR-Library estão incluídos em
\texttt{instancias/}, de modo que a execução não depende de conexão
com a Internet. Assim como no TP01, os scripts da raiz apenas compilam
e executam os testes automatizados:
\begin{lstlisting}
sh executar_testes.sh
# ou, no Windows:
executar_testes.bat
\end{lstlisting}
Para repetir a campanha oficial após a compilação, execute:
\begin{lstlisting}
java -cp out grafos.Main bench --dir instancias \
  --csv resultados/orlib.csv --limite-ms 3000
\end{lstlisting}
Os raios e tempos medidos estão em \path{resultados/orlib.csv}.
A lista de centros encontra-se em \path{resultados/centros.txt},
e os testes sintéticos em \path{resultados/sinteticos.csv}.
O relatório está disponível em PDF e em fonte TeX independente, que pode
ser compilada no diretório \texttt{relatorio/} com duas passagens do
\texttt{pdflatex}. O registro da auditoria independente acompanha os
demais resultados, sem exigir scripts adicionais para a entrega.

\section{Conclusão}
Foram implementadas e avaliadas duas abordagens para o problema dos
$k$-centros, atendendo às 40 entradas oficiais da atividade.
A implementação aproximada obteve gap médio de 32,69\% e ofereceu
soluções viáveis em todos os casos. A busca exata reduziu o raio em
33 casos, comprovou 13 ótimos e interrompeu 27
execuções pelo limite de tempo, mantendo os melhores centros encontrados.
Os arquivos-fonte, dados, tempos, raios, testes e PDF/TeX acompanham a
entrega, com procedimentos de reprodução e verificação independentes.

\begin{thebibliography}{9}
\bibitem{beasley}
BEASLEY, J. E. A note on solving large p-median problems.
\textit{European Journal of Operational Research}, v. 21, p. 270--273, 1985.
\bibitem{orlib}
BEASLEY, J. E. OR-Library: p-median -- uncapacitated. Disponível em:
\url{https://people.brunel.ac.uk/~mastjjb/jeb/orlib/pmedinfo.html}.
Acesso em: 16 nov. 2026.
\bibitem{gonzalez}
GONZALEZ, T. F. Clustering to minimize the maximum intercluster distance.
\textit{Theoretical Computer Science}, v. 38, p. 293--306, 1985.
\bibitem{enunciado}
PUC MINAS. \textit{Teoria dos Grafos e Computabilidade: Trabalho Prático n.º 02}.
Enunciado disponibilizado pela disciplina, 2025.

\bibitem{espelho}
RAFFAELLO. P-Median: Problems/pmed1--pmed40. Repositório público no GitHub.
Disponível em: \url{https://github.com/Raffaello/P-Median/tree/master/Problems}.
Acesso em: 16 nov. 2026.
\bibitem{yurtkuran}
YURTKURAN, A.; EMEL, E. A Modified Artificial Bee Colony Algorithm for
p-Center Problems. \textit{The Scientific World Journal}, 2014.
Disponível em: \url{https://pmc.ncbi.nlm.nih.gov/articles/PMC3926279/}.
\end{thebibliography}

\end{document}
